\section{\filetotitle}\label{sec:fewest_switches}
The \acrfull{fssh} method, developed by Tully in 1990,\supercite{10.1063/1.459170} is one of the most widely used \acrshort{tsh} algorithms. In \acrshort{fssh}, an ensemble of classical trajectories is propagated on adiabatic \acrshort{pes} according to Eq.~\eqref{eq:tsh_eom}, with stochastic hops between electronic states determined by the instantaneous electronic populations. The key idea is to minimize the number of hops while ensuring that the ensemble-averaged populations match the quantum mechanical populations.

Before we can obtain the hopping probabilities, it is convenient to express the Eq.~\eqref{eq:tsh_eom} in terms of the density matrix elements $\rho_{IJ}=c_I c_J^\ast$ as
\begin{equation}\label{eq:fssh_density_matrix_eom}
\i\dot\rho_{IJ}=\sum_K\left[\left(H_{IK}^{\symrm{el}}-\i\sigma_{IK}\right)\rho_{KJ}-\left(H_{KJ}^{\symrm{el}}-\i\sigma_{KJ}\right)\rho_{IK}\right],
\end{equation}
where the diagonal elements $\rho_{II}=\lvert c_I\rvert^2$ represent the populations of the adiabatic states, and the off-diagonal elements $\rho_{IJ}$ ($I\neq J$) represent the coherences between states. It is also useful to note that th electronic Hamiltonian (potential) is Hermitian and the \acrshort{tdc} is anti-Hermitian. From Eq.~\eqref{eq:fssh_density_matrix_eom}, we can derive the time evolution of the populations $\rho_{II}$ using the properties of complex numbers as
\begin{equation}\label{eq:fssh_population_eom}
\dot\rho_{II}=\sum_K b_{KI},
\end{equation}
where
\begin{equation}\label{eq:fssh_b_matrix}
b_{KI}=2\Im\left(H_{IK}^{\symrm{el}}\rho_{IK}^\ast\right)+2\Re\left(\sigma_{IK}\rho_{IK}^\ast\right).
\end{equation}
If we consider an adiabatic basis where the electronic Hamiltonian is diagonal, the first term in Eq.~\eqref{eq:fssh_b_matrix} vanishes, and the population transfer between states is mediated entirely by the \acrshort{tdc} term. We should also note that $b_{II}=0$ due to the properties of \acrshort{tdc} being anti-Hermitian and the density matrix being Hermitian.

Our goal is to derive hopping algorithm that minimizes the number of switches between states while ensuring that the ensemble-averaged populations match the quantum mechanical populations. Let's consider a two-state system for simplicity. In this case, the total number of trajectories at a time $t$ on state 1 be expressed as $N_1(t)=N\rho_{11}(t)$, where $N$ is the total number of trajectories in the ensemble. Without the loss of generality, we can assume that the population of state 1 is increasing, i.e., $\rho_{11}(t+\Delta t)>\rho_{11}(t)$. This means that the amount of trajectories that need to hop from state 2 to state 1 during the time interval $\Delta t$ must be larger than the number of trajectories hopping from state 1 to state 2. Since we want to minimize the number of hops, we can assume that no trajectories hop from state 1 to state 2 during this time interval. Therefore, the number of trajectories hopping from state 2 to state 1 must be equal to $N\left(\rho_{11}(t+\Delta t)-\rho_{11}(t)\right)$. The probability of a trajectory on state 2 hopping to state 1 during the time interval $\Delta t$ can then be expressed as the ratio of the number of trajectories hopping to the total number of trajectories on state 2 as
\begin{equation}\label{eq:fssh_hopping_probability_two_state}
P_{2\to 1}=\frac{N\left(\rho_{11}(t+\Delta t)-\rho_{11}(t)\right)}{N\rho_{22}}=\frac{\dot\rho_{11}\Delta t}{\rho_{22}}=\frac{b_{12}\Delta t}{\rho_{22}},
\end{equation}
where we used $\dot\rho_{11}=b_{12}$ from Eq.~\eqref{eq:fssh_population_eom}. The hopping probability from state 1 to state 2 can be derived similarly assuming that the population of state 1 is decreasing during the time interval $\Delta t$.

Should we want to generalize the hopping probabilities to more than two states, we can follow a similar reasoning. Let's assume that the population of state $I$ is increasing during the time interval $\Delta t$. The number of trajectories hopping from all other states to state $I$ must then be equal to $N\left(\rho_{II}(t+\Delta t)-\rho_{II}(t)\right)$. To minimize the number of hops, we can assume that no trajectories hop from state $I$ to any other state during this time interval. Therefore, the probability of a trajectory on state $J$ hopping to state $I$ during the time interval $\Delta t$ can be expressed as
\begin{equation}\label{eq:fssh_hopping_probability_multistate}
P_{J\to I}=\frac{b_{IJ}\Delta t}{\rho_{JJ}},
\end{equation}
where in the numerator we have the population flux only from state $J$ to state $I$, ignoring contributions from other states.
\subsection{Decoherence Correction}
In the \acrshort{fssh} method, we have two different ways to calculate the population of the electronic states. One is the usual ensemble-averaged population obtained by counting the number of trajectories on each state, and the other is the quantum mechanical population obtained from the density matrix elements $\rho_{II}$. One would expect these two populations to be equal, which is a condition known as internal consistency. However, in practice, the \acrshort{fssh} method often fails to maintain this internal consistency, meaning that the ensemble-averaged population does not match the quantum mechanical population.\supercite{10.1063/1.474288,10.1063/1.1793991,10.1063/1.471791}

In a truly quantum mechanical picture, when a wavepacket passes through a nonadiabatic coupling region, it can split into multiple components corresponding to different electronic states. These components usually end up evolving independently as they travel into distant regions of the \acrshort{pes}. The truly quantum mechanical time evolution therefore leads to a loss of coherence between the different components of the wavepacket.

Conversely, the \acrshort{fssh} method treats the system as a single point in phase space, which can only be on one electronic state at a time. When a trajectory hops from one state to another, it does not account for the fact that the wavepacket has split and that there should be a loss of coherence between the different components. This leads to an overestimation of the coherence between states, which in turn causes the ensemble-averaged population to deviate from the quantum mechanical population.

To address this issue, various decoherence correction schemes have been proposed. These schemes typically involve introducing a decoherence time scale or a decoherence factor that reduces the coherence between states after a hop occurs.\supercite{10.1063/1.1793991,10.1021/acs.jctc.2c00988,10.1021/jp991602b} We will describe here a scheme proposed by Granucci and Persico in 2007.\supercite{10.1063/1.2715585}

In this scheme, the electronic coefficients are modified after each integration step to account for the loss of coherence. Suppose the trajectory is currently on state $I$ and we are propagating the electronic coefficients from time $t$ to $t+\Delta t$. After the coefficient integration, we apply the transformation
\begin{equation}\label{eq:granucci_persico_decoherence}
c^\prime_J(t+\Delta t)=c_J(t+\Delta t)\exp\left(-\frac{\Delta t}{\tau_{IJ}}\right),\quad J\neq I,
\end{equation}
where $\tau_{IJ}$ is the decoherence time scale between states $I$ and $J$ and can be calculated as
\begin{equation}\label{eq:granucci_persico_decoherence_time_scale}
\tau_{IJ}=\frac{1}{\lvert E_I-E_J\rvert}\left(1+\frac{C}{E_{\symrm{kin}}}\right),
\end{equation}
where $E_I$ and $E_J$ are the energies of states $I$ and $J$, $E_{\symrm{kin}}$ is the kinetic energy of the particle, and $C$ is an empirical parameter. To ensure that the total population remains normalized after the transformation, we have to renormalize the current state coefficient as
\begin{equation}\label{eq:granucci_persico_decoherence_renormalization}
c^\prime_I(t+\Delta t)=C_I\sqrt{\frac{1-\sum_{J\neq I}\lvert c^\prime_J(t+\Delta t)\rvert^2}{\lvert c_I\rvert^2}}.
\end{equation}
Using this decoherence correction, the coherence between states is reduced after each integration step, which helps to maintain the internal consistency of the \acrshort{fssh} method and improve the agreement with quantum mechanical results. However, it is important to note that the choice of the empirical parameter $C$ can significantly affect the results, but we often choose $C=0.1$~Ha as it has been shown to provide good agreement with quantum mechanical results in many cases.\supercite{10.1063/1.1793991}
\subsection{Momentum Rescaling and Frustrated Hops}
Because different electronic states correspond to different potential energies at a given nuclear configuration, a hop between state $J$ and state $I$ results in a discontinuity in potential energy. To conserve the total energy of the system, the nuclear kinetic energy must be rescaled. It has been demonstrated that this energy exchange between the electronic and nuclear degrees of freedom should occur exclusively along the direction of the nonadiabatic coupling vector $\symbf{d}_{IJ}(\symbf{R})$. The momentum is adjusted such that
\begin{equation}\label{eq:fssh_energy_conservation}
\frac{\tilde{p}_I^2}{2M}+E_I(\symbf{R})=\frac{\tilde{p}_J^2}{2M}+E_J(\symbf{R}),
\end{equation}
where $\tilde{p}$ is the component of the nuclear momentum parallel to $\symbf{d}_{IJ}(\symbf{R})$, $M$ is the nuclear mass, and $\tilde{p}_J$ and $\tilde{p}_I$ are the momentum components along the coupling vector before and after the hop, respectively. If the algorithm dictates a hop to a higher energy state ($E_I>E_J$), but the kinetic energy available along the coupling vector is insufficient to overcome the energy gap (i.e., $\frac{\tilde{p}_J^2}{2M}<E_I-E_J$), the hop is forbidden. This event is known as a \textit{frustrated hop}.

To maintain balance within the \acrshort{fssh} ensemble, Tully proposed that during a frustrated hop, the electronic state remains unchanged, but the nuclear momentum component along the coupling vector is reversed ($\tilde{p}\to-\tilde{p}$), simulating an elastic bounce against the upper \acrshort{pes}. While empirically successful, this momentum reversal is an \textit{ad-hoc} algorithmic prescription rather than a consequence of the equations of motion.
